% !TeX root = ../main.tex
% ============================================================
% Chapter 5: Random Dungeon Map
% 说明：本文件由 main.tex 合并进主文档。章节编号 LaTeX 自动生成。
%       内容忠实来自 策划案.docx 原文（英文）。
%       图片：figures/Chapter05_RandomDungeonMap/
%       % TODO: 标记的小节是原文只有标题、没有内容的部分（待补充索引）。
% ============================================================
\section{Random Dungeon Map}

The random dungeon map is the biggest source of uncertainty and the reason players keep coming back to the same dungeon level.

\subsection{Dungeon Map Randomization Principles}

A random dungeon map should follow these principles:

\begin{enumerate}
  \item The road between the exit portal and the spawn point must be accessible.
  \item The terrain must be playable: it needs walls, plenty of open space for walking and fighting, and random obstacles that force players to use strategy.
\end{enumerate}

\subsection{How to Generate a Proper Dungeon Map?}

To generate a proper map, I recommend the steps below --- I have tried them, and they work (The new method is complete and working). This is the design I am most proud of, because it directly supports my core gameplay loop.

\textbf{New Method:}

\begin{enumerate}
  \item Use the cellular automaton algorithm to dig the terrain.
  \item Keep the biggest open area generated by the algorithm. (This is the core optimization step.)
  \item Place the Spawn Point and the Exit Portal randomly inside the open area. Put the spawn point in an open spot, and put the exit as far away as possible, also in an open spot.
  \item Scatter obstacles around the map.
  \item Check that the road between the exit portal and the spawn point stays accessible. If obstacles block the main road, scatter fewer obstacles.
  \item Make sure the Spawn Point and the Exit Portal sit on ground tiles; if not, re-run step 3.
  \item Tile the map and place the prefabs.
\end{enumerate}

Here Shows the effects of the algorithm.

\begin{figure}[htbp]
\centering
\includegraphics[width=0.85\textwidth]{dungeon_new_method.png}
\caption{New dungeon generation method}
\end{figure}

\textbf{Old Method:}

\begin{enumerate}
  \item Initialize the elements. Set the map scale first: the map is a square with a fixed border. Then place a spawn point and an exit portal in the right spots. In my design, the distance between them should be at least 80\% of the border length, and the exit portal should sit within 95\% of the border length so it never spawns on the border itself.
  \item Guarantee that the road between the exit portal and the spawn point is accessible. The A* algorithm can guarantee a spine road, but it needs optimization --- if we generate the spine road with A* first, the road has no randomization.
  \item Cellular automaton is a good fit for terrain generation. Use it at this step to generate random terrain. Its flaw: it tends to create karst-like terrain with stalactites, and the stalactites make wall generation awkward --- walls look strange when a stalactite has a pointed top.
  \item Cellular automaton may cover the spine road and break the path between the spawn point and the exit portal. This step needs optimization --- it becomes unnecessary when the earlier steps are handled properly.
  \item Once the spine road is generated, use BFS to inject more uncertainty into the map, giving players more adventurous decisions. Because of BFS's nature, it creates many side roads unrelated to the spine road, which can confuse the player.
  \item Dig rooms for the spawn point, the exit portal, combat, and chests. But rooms shouldn't sit right beside the spawn or exit points --- that would reduce the adventure uncertainty.
  \item Scatter random obstacles along the accessible roads. This step has no trick to it, but don't block the essential path with obstacles.
  \item Run an accessibility check --- the road must stay accessible. But this step is redundant if checking is already done in steps 2 and 4. To optimize: check the road once, and prepare a fallback plan in case the road is blocked.
  \item To make wall generation more reliable, cut off the stalactite's pointed top --- otherwise the walls cannot connect properly.
  \item Build the topology of the dungeon map and tile the map. Define each map element as an inside dot, a border dot, or an outside dot. This helps scripts understand tile locations: inside dots render floor, border dots render walls, outside dots render nothing.
  \item Create the exit portal prefab and the enemy spawn area. The enemy spawn area is described in section 5.3.
  \item Send messages to the state machine and the player. When the state machine receives the message, it changes its status to `dungeon'; the player is placed at the spawn point.
  \item Add Collider2D to walls and obstacles.
\end{enumerate}

Tips: the steps above have several defects, so I will reorder them to optimize the pipeline --- using fewer compute resources and generating a more playable dungeon map. The picture below shows version 1.0 of the random dungeon map.

\begin{figure}[htbp]
\centering
\includegraphics[width=0.85\textwidth]{dungeon_v1_defects.png}
\caption{V1.0 dungeon defects}
\end{figure}

Description of the picture elements:

\begin{enumerate}
  \item Player: the player's spawn point.
  \item Exit: the exit portal's location
  \item Defect 1: the walls don't connect properly, so the player can walk out of the dungeon.
  \item Defect 2: walls appear in the wrong places, making the map look strange.
  \item Defect 3: the exit's location is fixed, so the map doesn't show its randomness.
  \item Defect 4: the open areas are too small, so the player can't fight comfortably in the map.
\end{enumerate}

\subsection{Enemy/Enemy Difficulty Randomize}
% TODO: 原文仅有标题，内容待补充

\subsection{Treasure Randomize}
% TODO: 原文仅有标题，内容待补充
